
\documentclass{article}
\usepackage{colortbl}
\usepackage{makecell}
\usepackage{multirow}
\usepackage{supertabular}

\begin{document}

\newcounter{utterance}

\centering \large Interaction Transcript for game `cladder', experiment `full\_v1.5\_default', episode 4119 with reflexion{-}llama{-}3.1{-}8b{-}instruct.
\vspace{24pt}

{ \footnotesize  \setcounter{utterance}{1}
\setlength{\tabcolsep}{0pt}
\begin{supertabular}{c@{$\;$}|p{.15\linewidth}@{}p{.15\linewidth}p{.15\linewidth}p{.15\linewidth}p{.15\linewidth}p{.15\linewidth}}
    \# & \multicolumn{2}{c}{Player} && \multicolumn{2}{c}{Game Master} \\
    \hline

    \theutterance \stepcounter{utterance}  
    & & & \multicolumn{4}{p{0.6\linewidth}}{
        \cellcolor[rgb]{0.9,0.9,0.9}{
            \makecell[{{p{\linewidth}}}]{
                \texttt{\tiny{[P1$\langle$GM]}}
                \texttt{You are an expert in causal inference. The following question is not a typical commonsense query, but rather a meticulously designed question created by a professor specializing in causal inference, intended to assess the students' mastery of the course content.} \\
\texttt{Imagine a self{-}contained, hypothetical world with only the following conditions, and without any unmentioned factors or causal relationships: Eating citrus has a direct effect on vitmain C. Vitmain C has a direct effect on curly hair. For patients not consuming citrus, the probability of curly hair is 47\%. For patients consuming citrus, the probability of curly hair is 15\%. Will citrus intake decrease the chance of curly hair?} \\
\texttt{Start your answer with 'Yes' or 'No', followed by additional reasoning or evidence to support your explanation.} \\
            }
        }
    }
     \\ \\

    \theutterance \stepcounter{utterance}  
    & \multicolumn{4}{p{0.6\linewidth}}{
        \cellcolor[rgb]{0.561,0.737,0.561}{
            \makecell[{{p{\linewidth}}}]{
                \texttt{\tiny{[P1$\rangle$GM]}}
                \texttt{No.} \\
\\ 
\texttt{To answer this question, we need to apply the concept of causal inference, specifically the backdoor criterion or the frontdoor adjustment. } \\
\\ 
\texttt{First, let's establish the causal relationships: } \\
\\ 
\texttt{{-} Eating citrus (C) causes an increase in vitamin C (V).} \\
\texttt{{-} Vitamin C (V) causes curly hair (H).} \\
\\ 
\texttt{We are given the following probabilities:} \\
\texttt{{-} P(H|not C) = 0.47 (probability of curly hair for patients not consuming citrus)} \\
\texttt{{-} P(H|C) = 0.15 (probability of curly hair for patients consuming citrus)} \\
\\ 
\texttt{The question asks whether citrus intake decreases the chance of curly hair. To answer this, we need to compare the probability of curly hair when not consuming citrus (P(H|not C)) to the probability of curly hair when consuming citrus (P(H|C)).} \\
\\ 
\texttt{However, the given probabilities do not directly answer the question. Instead, we need to consider the causal relationship between citrus intake and curly hair. } \\
\\ 
\texttt{To determine whether citrus intake decreases the chance of curly hair, we need to consider the effect of vitamin C on curly hair. Since vitamin C causes curly hair, and citrus intake increases vitamin C, we can infer that citrus intake indirectly causes curly hair.} \\
\\ 
\texttt{However, the question asks whether citrus intake decreases the chance of curly hair. To determine this, we need to consider the effect of citrus intake on the probability of curly hair, while accounting for the effect of vitamin C.} \\
\\ 
\texttt{Using the backdoor criterion, we can calculate the causal effect of citrus intake on curly hair as follows:} \\
\\ 
\texttt{{-} P(H|C) = P(H|not C) + P(V|C) * (P(H|V) {-} P(H|not V))} \\
\\ 
\texttt{where P(V|C) is the probability of vitamin C given citrus intake, and P(H|V) and P(H|not V) are the probabilities of curly hair given vitamin C and not vitamin C, respectively.} \\
\\ 
\texttt{However, we are not given the exact values of P(V|C), P(H|V), and P(H|not V). Nevertheless, we can still reason about the causal effect of citrus intake on curly hair.} \\
\\ 
\texttt{Since vitamin C causes curly hair, and citrus intake increases vitamin C, we can infer that citrus intake indirectly causes curly hair. However, the question asks whether citrus intake decreases the chance of curly hair.} \\
\\ 
\texttt{To determine this, we need to consider the effect of citrus intake on the probability of curly hair, while accounting for the effect of vitamin C. } \\
\\ 
\texttt{Given that the probability of curly hair for patients consuming citrus (P(H|C)) is lower than the probability of curly hair for patients not consuming citrus (P(H|not C)), it may seem that citrus intake decreases the chance of curly hair. However, this conclusion is based on a naive comparison of the probabilities.} \\
\\ 
\texttt{To accurately determine the causal effect of citrus intake on curly hair, we need to consider the effect of vitamin C on curly hair. Since vitamin C causes} \\
            }
        }
    }
    & & \\ \\

    \theutterance \stepcounter{utterance}  
    & & & \multicolumn{2}{p{0.3\linewidth}}{
        \cellcolor[rgb]{0.9,0.9,0.9}{
            \makecell[{{p{\linewidth}}}]{
                \texttt{\tiny{[GM$|$GM]}}
                \texttt{No} \\
            }
        }
    }
    & & \\ \\

    \theutterance \stepcounter{utterance}  
    & & & \multicolumn{2}{p{0.3\linewidth}}{
        \cellcolor[rgb]{0.9,0.9,0.9}{
            \makecell[{{p{\linewidth}}}]{
                \texttt{\tiny{[GM$|$GM]}}
                \texttt{yes} \\
            }
        }
    }
    & & \\ \\

    \theutterance \stepcounter{utterance}  
    & & & \multicolumn{2}{p{0.3\linewidth}}{
        \cellcolor[rgb]{0.9,0.9,0.9}{
            \makecell[{{p{\linewidth}}}]{
                \texttt{\tiny{[GM$|$GM]}}
                \texttt{game\_result = LOSE} \\
            }
        }
    }
    & & \\ \\

    \theutterance \stepcounter{utterance}  
    & & \multicolumn{4}{p{0.6\linewidth}}{
        \cellcolor[rgb]{0.95,0.95,0.95}{
            \makecell[{{p{\linewidth}}}]{
                \texttt{\tiny{[P1$\leftrightarrow$P]}}
                \texttt{Post{-}episode Reflexion (not sent to GM):} \\
\texttt{New plan: } \\
\\ 
\texttt{In my previous attempt, I failed to accurately determine the causal effect of citrus intake on curly hair. I correctly identified the causal relationships between eating citrus, vitamin C, and curly hair, but I failed to properly apply the backdoor criterion to calculate the causal effect of citrus intake on curly hair.} \\
\\ 
\texttt{To improve my strategy, I will focus on correctly applying the backdoor criterion to calculate the causal effect of citrus intake on curly hair. Specifically, I will make sure to:} \\
\\ 
\texttt{1. Identify the correct causal relationships between eating citrus, vitamin C, and curly hair.} \\
\texttt{2. Calculate the probability of vitamin C given citrus intake (P(V|C)).} \\
\texttt{3. Calculate the probability of curly hair given vitamin C (P(H|V)) and not vitamin C (P(H|not V)).} \\
\texttt{4. Use the backdoor criterion to calculate the causal effect of citrus intake on curly hair.} \\
\\ 
\texttt{To do this, I will need to carefully analyze the given probabilities and causal relationships, and make sure to apply the backdoor criterion correctly. I will also need to be careful not to make any assumptions or naive comparisons that may lead to incorrect conclusions.} \\
\\ 
\texttt{By following this strategy, I hope to accurately determine the causal effect of citrus intake on curly hair and provide a correct answer to the question.} \\
            }
        }
    }
    & \\ \\

\end{supertabular}
}

\end{document}
